\subsection{\identifier{UC-05} --- Manage Properties and Parcels}
\label{uc:05}

\paragraph{Purpose}

Allow an authorised user to create, consult, update,
archive, reactivate and, where permitted, delete
properties and their optional parcels.

The current agricultural scope concerns
olive harvesting.


\paragraph{Actors}

\textbf{Primary actor:} Farm Owner.

\textbf{Additional actor:} Field Manager,
when explicitly granted property-management
permissions by the owner.

Permission to record operational work does not
automatically grant permission to manage properties.

\paragraph{Preconditions}

\begin{itemize}
    \item The user has permission for
    the requested operation.
    \item Creation, modification, archival,
    reactivation and deletion require
    an internet connection.
    \item For operations on an existing record,
    that record exists within the authorised scope.
    \item To create a parcel, its property exists.
    \item Offline consultation requires
    previously downloaded authorised information.
\end{itemize}

\paragraph{Main Flow}

\begin{enumerate}
    \item The user opens property management.
    \item The user chooses to create a property
    or selects an existing property.
    \item The system presents the existing
    information or an entry form.
    \item The user enters or updates the property
    name and textual location.
    \item The user may enter the declared area
    in hectares and one or more olive varieties.
    \item The system identifies the owner
    associated with the property.
    \item If internal divisions are needed,
    the user creates or updates named parcels.
    Parcel area and olive varieties are optional.
    \item The system validates the information
    and access permissions.
    \item The user reviews and confirms the changes.
    \item The system saves the records
    and confirms the result.
\end{enumerate}

For consultation only, the flow ends after
the requested information is displayed.

\paragraph{Alternative Flows and Exceptions}

\begin{description}

    \item[A1 --- No internal divisions are needed.]
    The user saves the property without parcels.
    The absence of parcels does not prevent
    use of the property in operational records.

    \item[A2 --- Optional information is unknown.]
    The user leaves area or variety information
    unfilled and may add it later.
    An unknown area is not treated as zero.

    \item[A3 --- The user archives a record.]
    The system requests confirmation and
    archives the selected property or parcel.
    Historical records remain accessible
    within the user's permissions.
    The archived record is unavailable
    for new operational associations.

    \item[A4 --- The user reactivates a record.]
    The system validates the request
    and restores its availability,
    subject to its parent property's state.

    \item[A5 --- The user requests deletion.]
    The system checks for associated records.
    If none exist, the user may confirm deletion.
    Otherwise, deletion is refused
    and archival is offered.

    \item[A6 --- Connectivity is unavailable.]
    The system permits consultation of
    previously downloaded authorised records
    and identifies their synchronisation status.
    Management operations are unavailable
    until connectivity is restored.
\end{description}

\paragraph{Postconditions}

\textbf{On successful creation or update:}
Property and parcel information is saved
with the correct ownership and parent associations.

\textbf{On archival:}
Historical information is preserved,
and new operational use is blocked
for the archived scope.

\textbf{On reactivation:}
The record becomes available again,
subject to the property's availability.

\textbf{On deletion:}
Only the eligible record without
associated records is removed.

\textbf{On consultation, cancellation or rejection:}
Stored information remains unchanged.

\paragraph{Related Requirements and Rules}
\hyperref[fr:009]{\identifier{FR-009}}, \hyperref[fr:010]{\identifier{FR-010}}, \hyperref[fr:011]{\identifier{FR-011}}, \hyperref[fr:012]{\identifier{FR-012}}, \hyperref[fr:013]{\identifier{FR-013}}.

\hyperref[br:005]{\identifier{BR-005}}, \hyperref[br:006]{\identifier{BR-006}}, \hyperref[br:013]{\identifier{BR-013}}.

\paragraph{Decision References}
\decisionref{Q-SCOPE-01}, \decisionref{D-06}.
